% Methods and Results sections, IEEEtran (pdfLaTeX).
% Preamble requirements: \usepackage{amsmath,amssymb,graphicx,booktabs,siunitx}
% Bibliography keys used: apriltag, px4, mavsdk.
% Lines marked "AUTHOR QUERY" must be resolved and stripped before submission.

\section{Methods}

We compare two landing controllers that differ only in how the horizontal velocity
command is generated. The \emph{vision-guided controller} servoes the vehicle on the
image-plane position of a fiducial marker mounted on the platform deck. The
\emph{barometric controller} descends at a constant rate with no horizontal command and
serves as the baseline. Both controllers stream velocity setpoints to the flight
controller at \SI{50}{\hertz} through the vehicle's onboard autopilot
interface\footnote{PX4 v1.14.}~\cite{px4,mavsdk}. Fig.~\ref{fig:block} shows the
arrangement.

\begin{figure}[!t]
  \centering
  \includegraphics[width=\columnwidth]{fig_block_diagram.pdf}
  \caption{Landing controller architecture. Fiducial detections from the downward camera
  are smoothed and differenced against the image center to form the horizontal velocity
  command, while the fused altitude estimate gates descent and selects the touchdown
  rate. The barometric controller omits the camera, detector, and image-error blocks.}
  \label{fig:block}
\end{figure}

\subsection{Landing Guidance from Image-Plane Error}

A downward-facing camera observes a square fiducial marker at the center of the platform
deck, and a detector reports the marker center in pixel coordinates~\cite{apriltag}.
Because the camera looks along the vehicle's vertical axis, the two image axes map
one-to-one onto the two horizontal body axes, and driving the marker to the image center
drives the vehicle over the deck center. No range or pose reconstruction is therefore
required, and the guidance law acts directly on pixel error.

Detections arrive an order of magnitude more slowly than setpoints are streamed, so most
control cycles execute without a new measurement. Two mechanisms bridge the gap. The
detected marker center is smoothed by an exponential moving average filter, and between
detections the controller continues to act on the most recent smoothed estimate. Let
$(\bar{u},\bar{v})$ denote the smoothed marker center and $(u_0,v_0)$ the image center.
The normalized image error is given by
\begin{equation}
  e_u = \frac{\bar{u}-u_0}{u_0}, \qquad
  e_v = \frac{\bar{v}-v_0}{v_0},
  \label{eq:err}
\end{equation}
where $e_u$ is the lateral image error and $e_v$ is the longitudinal image error, each
normalized so that unity corresponds to the marker lying at the image border.

The horizontal velocity command is given by
\begin{equation}
  \begin{aligned}
    v_x &= \operatorname{sat}\!\left(-k_p\,e_v - k_d\,\dot{x},\; v_h^{\max}\right),\\
    v_y &= \operatorname{sat}\!\left(-k_p\,e_u - k_d\,\dot{y},\; v_h^{\max}\right),
  \end{aligned}
  \label{eq:guidance}
\end{equation}
where $v_x$ and $v_y$ are the commanded longitudinal and lateral velocities; $k_p$ is the
proportional gain on normalized image error; $k_d$ is the damping gain; $\dot{x}$ and
$\dot{y}$ are the measured horizontal velocities of the vehicle; $v_h^{\max}$ is the
horizontal speed limit; and $\operatorname{sat}(a,b)=\min(\max(a,-b),b)$ is the symmetric
saturation.

The damping term in~\eqref{eq:guidance} acts on the measured vehicle velocity rather than
on a numerical derivative of the image error. Differentiating the image error would
amplify detection noise at the camera frame rate, whereas the velocity estimate is
available at the full setpoint rate.

When no detection has arrived for longer than a fixed timeout, the controller applies a
halved copy of the last command it issued. This decay hold keeps the vehicle drifting
with the platform instead of stopping over ground during a dropout, and it is the
behavior invoked by the dropout events reported in Section~\ref{sec:results}.

\subsection{Altitude Estimation and Touchdown}

Two altitude sources are available. The barometer is continuous but drifts, and the
downward rangefinder is unbiased but noisy near the ground and intermittent. The altitude
estimate is given by
\begin{equation}
  \hat{h}_k = (1-\kappa)\,h_k^{\mathrm{baro}} + \kappa\,h_k^{\mathrm{range}},
  \qquad \kappa = \frac{\Delta t}{\tau + \Delta t},
  \label{eq:alt}
\end{equation}
where $\hat{h}_k$ is the altitude estimate at cycle $k$; $h_k^{\mathrm{baro}}$ and
$h_k^{\mathrm{range}}$ are the barometric and rangefinder measurements; $\Delta t$ is the
control period; and $\tau$ is the blending time constant. When no rangefinder return is
available, the barometric measurement is used directly. The estimate is thus a
first-order blend that low-pass filters the rangefinder against the barometer, and it
does not high-pass the barometric channel.
% AUTHOR QUERY: the implementation names this filter "complementary" but computes the
% blend above, in which the barometric term enters at full magnitude rather than as an
% increment. Equation (3) reports what is implemented. Confirm this is intended before
% submission, and correct either the code or this paragraph, not only the paragraph.

Descent is gated on alignment. The vertical velocity command is given by
\begin{equation}
  v_z =
  \begin{cases}
    -v_z^{\mathrm{td}}, & \hat{h} < h_f,\\[2pt]
    -v_z^{\max}, & \hat{h} \ge h_f \ \text{and}\ \lVert \bar{\mathbf{p}} - \mathbf{p}_0 \rVert < \rho,\\[2pt]
    0, & \text{otherwise},
  \end{cases}
  \label{eq:vz}
\end{equation}
where $v_z$ is the commanded vertical velocity, positive upward; $v_z^{\mathrm{td}}$ is
the reduced touchdown descent rate; $h_f$ is the flare altitude; $v_z^{\max}$ is the
cruise descent rate; $\bar{\mathbf{p}}$ and $\mathbf{p}_0$ are the smoothed marker center
and the image center; and $\rho$ is the alignment radius in pixels. Above the flare
altitude the vehicle descends only while the marker lies within $\rho$ of the image
center, so horizontal error is reduced before altitude is given up. Below the flare
altitude the first case of~\eqref{eq:vz} takes precedence, and the vehicle descends at
the reduced rate whether or not it is aligned. Committing to touchdown in this way avoids
holding altitude in the region where the rangefinder is noisiest and where ground effect
perturbs the horizontal loop, at the cost of leaving any residual alignment error
uncorrected over the final $h_f$ of descent.

\subsection{Baseline Controller}

The barometric controller commands zero horizontal velocity and a constant descent rate
$v_z^{\max}$ until touchdown. It neither gates descent on alignment nor reduces the
descent rate near the ground, and it does not read the altitude estimate at all, so it
represents open-loop descent onto the platform's last known position. Comparing against
it isolates the contribution of the image-plane guidance law, because the two
controllers share the vehicle, the setpoint interface, and the cruise descent rate.

\subsection{Experimental Setup}

Experiments used a \SI{450}{\milli\meter} quadrotor carrying a downward-facing
global-shutter camera at $640\times480$ pixels. The camera streams at \SI{8}{\hertz},
against a \SI{50}{\hertz} setpoint rate. The landing target was a
\SI{160}{\milli\meter} AprilTag
36h11 marker at the center of a cart deck of \SI{300}{\milli\meter} radius, towed at a
constant \SI{0.3}{\meter\per\second}. Fig.~\ref{fig:setup} shows the apparatus.

Ground truth came from an OptiTrack motion capture system sampling at
\SI{120}{\hertz}, which tracked both the vehicle and the deck throughout each descent.
Gains, filter constants, and the alignment radius were selected empirically for this
airframe and camera, and represent values tuned for this apparatus rather than
recommended settings.

We flew 12 landings with each controller, alternating between them so that each
vision-guided landing is paired with a barometric landing flown under matched platform
and wind conditions.
% AUTHOR QUERY: the lab notes record a paired analysis but do not state the pairing
% scheme. Confirm the alternation described here, or replace with the scheme actually
% flown.

\begin{figure}[!t]
  \centering
  \includegraphics[width=\columnwidth]{fig_setup.pdf}
  \caption{Experimental apparatus. Callouts identify the quadrotor, the downward camera,
  the AprilTag marker at the deck center, the \SI{300}{\milli\meter} deck boundary, and
  the motion capture markers used for ground truth.}
  \label{fig:setup}
\end{figure}

\subsection{Metrics and Statistical Analysis}

The primary metric is the touchdown offset, defined as the horizontal distance between
the vehicle center and the deck center at the instant of touchdown, measured by motion
capture.
% AUTHOR QUERY: confirm the reference points (vehicle center versus landing-gear
% centroid) used to compute the offset from the motion capture data.
A landing is recorded as contained when the touchdown offset is smaller than the
\SI{300}{\milli\meter} deck radius. Descent duration is the interval between the start of
the commanded descent and touchdown.

Touchdown offsets are summarized as mean $\pm$ standard deviation, which reports the
absolute accuracy and its repeatability across trials. The two controllers are compared
by a two-sided paired $t$-test on the touchdown offset, with the paired difference and its
\SI{95}{\percent} confidence interval reporting the size and precision of the change.
Because a single confirmatory comparison was specified, no correction for multiple
comparisons was applied. Containment and descent duration are reported descriptively.

\section{Results}
\label{sec:results}

Table~\ref{tab:results} summarizes both controllers over the 12 paired landings. Image
guidance reduced the mean touchdown offset by roughly three quarters, from
\SI{187}{\milli\meter} to \SI{46}{\milli\meter} (paired difference
\SI{141}{\milli\meter}, \SI{95}{\percent} CI [98, 184], $p = 2\times10^{-4}$).
Trial-to-trial dispersion fell by a factor of 5.1, from a standard deviation of
\SI{92}{\milli\meter} to \SI{18}{\milli\meter}.

All 12 vision-guided landings were contained within the deck radius. Three of the 12
barometric landings were not. Fig.~\ref{fig:offsets} shows the per-trial offsets against
the deck boundary.

Descent duration was $14.2 \pm \SI{1.9}{\second}$ for the vision-guided controller and
$9.8 \pm \SI{0.7}{\second}$ for the barometric controller, a difference of
\SI{4.4}{\second}.

In one vision-guided trial the marker was lost twice for longer than the detection
timeout. The decay hold was applied on both occasions, and detections resumed before
touchdown. No other vision-guided trial recorded a dropout beyond the timeout.

\begin{table}[!t]
  \centering
  \caption{Landing Performance Over 12 Paired Trials. Differences Are Vision-Guided
  Minus Barometric}
  \label{tab:results}
  \setlength{\tabcolsep}{3pt}
  \footnotesize
  \begin{tabular}{lccc}
    \toprule
    & Vision-guided & Barometric & Difference [95\% CI] \\
    \midrule
    Touchdown offset (mm) & $46 \pm 18$ & $187 \pm 92$ & $-141$ $[-184, -98]$ \\
    Within deck radius    & 12/12 & 9/12 & --- \\
    Descent duration (s)  & $14.2 \pm 1.9$ & $9.8 \pm 0.7$ & $+4.4$ \\
    \bottomrule
  \end{tabular}
\end{table}

\begin{figure}[!t]
  \centering
  \includegraphics[width=\columnwidth]{fig_offsets.pdf}
  \caption{Touchdown offsets for each of the 12 paired landings. Filled circles mark
  vision-guided landings and open squares mark barometric landings, with paired trials
  joined by a line. The horizontal dashed line marks the \SI{300}{\milli\meter} deck
  radius, beyond which a landing missed the platform.}
  \label{fig:offsets}
\end{figure}
